\chapter{Qubitization and Quantum Signal Processing}
\label{ch:lec08}

\begin{objectives}
\begin{itemize}[leftmargin=1.4em,nosep]
  \item Derive Chebyshev transformations from a $2\times2$ rotation.
  \item Qubitize Hermitian and general block encodings.
  \item State the quantum signal processing (QSP) theorem.
  \item Survey methods for finding phase factors.
\end{itemize}
\end{objectives}

\todo{Write the introductory paragraph orienting the lecture.}

\section{Chebyshev polynomials from a $2\times2$ block}
\label{sec:lec08-1}
\todo{Draft this section.}

\section{Qubitization of Hermitian matrices}
\label{sec:lec08-2}
\todo{Draft this section.}

\section{Qubitization of general matrices}
\label{sec:lec08-3}
\todo{Draft this section.}

\section{Quantum signal processing}
\label{sec:lec08-4}
\todo{Draft this section.}

\section{Finding phase factors}
\label{sec:lec08-5}
\todo{Draft this section.}

\exercisesection
\begin{exercise}
\todo{Add exercise.}
\end{exercise}
\begin{exercise}
\todo{Add exercise.}
\end{exercise}
\begin{exercise}
\todo{Add exercise.}
\end{exercise}

\readingsection
\begin{itemize}[leftmargin=1.4em]
  \item Lin Lin, \emph{qasc notes}, Ch.~7.
  \item Lin Lin, \emph{live notes}, Ch.~10 (qubitization) and Ch.~12 (QSP).
  \item Dong An, \emph{Quantum Numerical Linear Algebra} talk (qubitization, QSP).
\end{itemize}
